\documentclass{article}

\usepackage{tabularx}
\usepackage{booktabs}

\title{Problem Statement and Goals\\\progname}

\author{\authname}

\date{}

\input{../Comments.text}
\input{../Common.text}

\begin{document}

\maketitle

\begin{table}[hp]
\caption{Revision History} \label{TblRevisionHistory}
\begin{tabularx}{\textwidth}{llX}
\toprule
\textbf{Date} & \textbf{Developer(s)} & \textbf{Change}\\
\midrule
2026/09/25 & Jake Finlay & Completed section 1.4 \\
\bottomrule
\end{tabularx}
\end{table}

\section{Problem Statement}

\wss{You should check your problem statement with the
\href{https://smiths.github.io/capTemplate/Checklists/ProbState-Checklist.pdf}
{problem statement checklist}}. 

\wss{You can change the section headings, as long as you include the required
information.}

\subsection{Problem}

\wss{What problem are you trying to solve?  Why is it important? Do not talk
about how you will solve the problem, only what the problem is.  AI is rarely
part of the problem.  The problem is ``what'' you want to do, not ``how'' you
want to do it.  The problem statement should be written in a way that is
understandable to a non-technical audience.}

\subsection{Inputs and Outputs}

\wss{Characterize the problem in terms of ``high level'' inputs and outputs.  
Use abstraction so that you can avoid details.}

\subsection{Stakeholders}

\subsection{Environment}

\subsubsection{Software}

The application will run on iOS and Android operating systems. The application will be distributed through the Apple App Store and Google Play Store.

\subsubsection{Hardware}

The application will run on consumer smartphones and tablets. An internet connection will be required to download the application.

\section{Goals}

\section{Stretch Goals}

\wss{Goals you hope to get to, but not necessary.  They are also helpful if the
scope of the original project needs to be expanded.}

\section{Custom Documents (CDs)}

\wss{For CAS 741: State whether the project is a research project. This
designation, with the approval (or request) of the instructor, can be modified
over the course of the term.}

\wss{For SE Capstone: List your custom documents (CDs).  Potential CDs include
usability testing, code walkthroughs, user documentation, formal proof,
GenderMag personas, Design Thinking, etc.  (The full list is on the course
outline and in Lecture 02.) Normally the number of CDs will be two (one per
term).  Approval of the CDs will be part of the discussion with the instructor
for approving the project. The CDs, with the approval (or request) of the
instructor, can be modified over the course of the term.}

\newpage{}

\section*{Appendix --- Reflection}

\wss{Not required for CAS 741}

\input{../Reflection.text}

\begin{enumerate}
    \item What went well while writing this deliverable?
    \item What pain points did you experience during this deliverable, and how
    did you resolve them?
    \item How did you and your team adjust the scope of your goals to ensure
    they are suitable for a Capstone project (not overly ambitious but also of
    appropriate complexity for a senior design project)?
\end{enumerate}

\subsection*{Ethan McMehen}

\subsubsection*{What went well while writing this deliverable?}

I specifically wrote the stakeholders section of the document and I found that 
explicitly documenting the stakeholders and their relationships in the tier list 
really helped to identify who the application is being developed for. By writing 
the descriptions of each stakeholder and their relation to the problem statement, 
it really helped me to prepare for requirements elicitation and who to approach 
for what kind of information. I also found that writing a problem statement is 
fantastic too because it helps to focus on our specific problem we are aiming to 
solve. Our project description by our supervisor was solid, but writing our own 
version helps us to have a definition in an engineering context that we can refer 
to throughout the project. Overall, I learned that by having explicit definitions 
and descriptions for topics or information that seems relatively straight forward 
is still important to have a unified mission.

\subsubsection*{What pain points did you experience during this deliverable, and how did you resolve them?}

My pain points that I experienced during this deliverable have to be creating 
the right descriptions for the stakeholders and their relationship to the 
application that we are going to build for the capstone. While the stakeholders 
of children, parents, ophthalmologists etc. is easy to name, actually coming 
down to their effect and purpose in the project required some thinking. To 
resolve this pain point, I broke down their role into three parts: are they 
using the app, are they evaluating the app effectiveness, and do they care about 
the safety and ability for the app to reduce stress. I used these points to 
devise descriptions about each stakeholder relative to their involvment in each 
of these three parts. This helped me reinforce my learning that breaking down 
the problem or pain point into smaller pieces is a great approach for software 
documentation.

\end{document}